El \dfn{frontend} se construye sobre un \dfn{backend} que ya existía y que no forma parte
de este trabajo: el servicio Django REST del TFG de backend~\cite{scdee_backend_repo}.
La primera tarea fue, por tanto, de análisis y no de diseño: averiguar qué ofrece realmente
ese servicio, con qué reglas, y en qué se diferenciaba de lo que suponía el prototipo de
interfaz heredado. Esta sección describe cómo se hizo y qué se obtuvo; las conclusiones
alimentan los requisitos de las secciones siguientes.

\subsection{Punto de partida}

El repositorio del \dfn{frontend} contenía un prototipo del módulo del profesor (rama
\texttt{modulo-profesor}, 15 commits) desarrollado antes de que el \dfn{backend} estuviera
disponible. Su documento de contrato describía una primera fase del servicio en la que solo
existía la comprobación de salud y una lista de operaciones \emph{planificadas}. El
\dfn{backend} real se localizó en su propio repositorio, rama \texttt{backend}, cuyo último
commit (\texttt{a8eb9f43}, 19/05/2026) se tomó como versión de referencia para todo el
trabajo.

La obtención del código ya planteó un primer obstáculo: tanto el clonado completo como el
superficial fallaban al descargar ficheros binarios voluminosos de la documentación. Se
resolvió consultando primero el árbol del repositorio mediante la \acs{api} de GitHub (791
entradas) para decidir qué hacía falta, y clonando después de forma parcial
(\texttt{--filter=blob:none}) con \textit{sparse checkout} limitado al código del servicio y a
la especificación OpenAPI.

\subsection{El contrato como fuente de verdad}

El \dfn{backend} publica una descripción OpenAPI~3.0.3 generada con drf-spectacular
(\texttt{docs/files/spectacular/SCDEE API.yaml}, unos 534\,KB)~\cite{codex_openapi303,codex_drfspectacular}.
Contiene 95 operaciones sobre 73 rutas ---94 si se descuenta la comprobación de salud---,
127 esquemas de datos y un único esquema de seguridad, \textit{Bearer} con \acs{jwt}.

Un documento de ese tamaño no se puede consultar de forma fiable a mano ni, por la misma
razón, dejárselo leer entero a un \dfn{agente} cada vez que necesita un dato. Se escribieron
dos \textit{scripts} deterministas (\texttt{scripts/analyze/endpoints.py} y
\texttt{schemas.py}) que extraen del YAML una línea por operación (método, ruta, etiqueta,
cuerpo de la petición, respuestas y resumen) y una ficha por esquema (campos, tipos,
obligatoriedad). El resultado, \texttt{docs/frontend/API\_INVENTORY.md}, se convirtió en el
índice de consulta del proyecto y en una regla explícita de trabajo: ninguna pantalla usa
una operación o un campo que no figure en él (\cref{AP:API}). El mismo inventario sirve
después para medir la cobertura del contrato (\cref{SEC:TEST-COVERAGE}).

La extracción automática resolvió la parte estructural: operaciones, formas de los mensajes,
paginación por desplazamiento (\texttt{count}, \texttt{page}, \texttt{page\_size},
\texttt{total\_pages}, \texttt{results}), formato simbólico de los errores
(\texttt{error\_code} y, en validación, códigos por campo) y nombres exactos de los
\textit{tokens} (\texttt{access\_token}, \texttt{refresh\_token}).

\subsection{Lo que el esquema no expresa}

Un esquema describe la forma de los mensajes, pero no quién puede ver qué, en qué momento ni
qué ocurre si dos personas actúan a la vez. Para una interfaz esas reglas son tan
importantes como los campos, porque deciden qué botones se muestran, qué variantes tiene
cada vista y qué errores hay que explicar. Se leyeron por ello las partes del código del
\dfn{backend} donde el esquema podía no reflejar el comportamiento:
\begin{description}
  \item[Identidad y roles.] El complemento de autenticación (\texttt{jwt\_plugin.py}) añade
        al \acs{jwt} los indicadores \texttt{is\_staff} e \texttt{is\_superadmin} y la
        organización; la respuesta del inicio de sesión no los incluye. El rol docente no es
        global: \texttt{MySubjectsView} devuelve un rol por asignatura
        (\cref{SEC:AN-ACTORS}).
  \item[Permisos.] La tabla \texttt{DEFAULT\_PERMISSIONS} fija once permisos
        \texttt{can\_*} por rol, que pueden sobrescribirse por membresía.
  \item[Ciclo de vida de una \dfn{instancia}.] \texttt{VALID\_TRANSITIONS} define una
        máquina de estados de diez estados: \texttt{ASSEMBLING} $\rightarrow$
        \texttt{RECEIVED} $\rightarrow$ \texttt{QUEUED} $\rightarrow$
        \texttt{PENDING\_GRADING} $\rightarrow$ \texttt{GRADED} $\rightarrow$
        \texttt{PUBLISHED} $\rightarrow$ \texttt{IN\_REVIEW} $\rightarrow$
        \texttt{FINALIZED} $\rightarrow$ \texttt{ARCHIVED}, con un estado
        \texttt{PENDING\_REVIEW} para las revisiones aceptadas y varios retrocesos
        permitidos (por ejemplo, de \texttt{GRADED} a \texttt{PENDING\_GRADING}).
  \item[Visibilidad para el estudiante.] \texttt{access\_control.py} solo muestra la
        información básica a partir de la publicación; el acceso a páginas, anotaciones,
        \dfn{rubrica} o desglose depende además de seis permisos del examen
        (\texttt{student\_permissions}) y de que el periodo de revisión esté abierto, salvo
        que el examen permita seguir viendo las páginas tras la revisión.
  \item[Concurrencia.] \texttt{grading.py} aplica control optimista: para modificar una
        nota existente hay que enviar la anterior (\texttt{old\_score}); si no coincide, el
        servicio responde 409 con \texttt{VERSION\_CONFLICT}.
  \item[Representación según el rol.] Un mismo recurso se serializa de forma distinta según
        quién lo pide (por ejemplo, tres representaciones del examen y una específica de la
        \dfn{instancia} para el estudiante), de modo que la interfaz no puede suponer los
        mismos campos para todos.
  \item[Anotaciones.] El contenido de una anotación se guarda en almacenamiento de objetos
        y la respuesta solo ofrece una dirección de descarga firmada
        (\texttt{download\_url}), no el contenido.
\end{description}
Algunas de estas lecturas sirvieron también para descartar hipótesis. Al revisar
\texttt{MySubjectsView} se sospechó una fuga entre organizaciones, porque busca el curso
activo sin filtrar explícitamente por organización; la lectura del gestor
\textit{multi-tenant} del modelo mostró que el filtro se aplica siempre, así que no hay tal
fuga. Se documenta como ejemplo de una sospecha que se verificó antes de comunicarla.

\subsection{Desajustes del prototipo previo}

La comparación del prototipo con el \dfn{backend} real arrojó seis desajustes:
\begin{enumerate}
  \item El inicio de sesión leía \texttt{access}/\texttt{refresh}, pero el servicio devuelve
        \texttt{access\_token}/\texttt{refresh\_token}; la renovación y el cierre de sesión
        enviaban el campo con el nombre equivocado.
  \item El rol se deducía de \texttt{is\_staff} como si fuera global (profesor o estudiante),
        cuando \texttt{is\_staff} identifica al gestor de la organización y el rol docente
        depende de la asignatura.
  \item Se esperaba \texttt{is\_staff} en la respuesta del inicio de sesión, pero solo
        viaja dentro del \acs{jwt}.
  \item La precarga de anotaciones leía un campo \texttt{payload} que la respuesta no
        contiene.
  \item Se usaban estados simplificados de la \dfn{instancia} que no existen en el servicio.
  \item El estado de la corrección y las anotaciones se guardaban en
        \texttt{localStorage} como fuente de verdad.
\end{enumerate}
Afectaban a la autenticación, al modelo de roles y a la persistencia, es decir, a la base de
toda la aplicación. Se concluyó que era más económico reescribir la capa de datos y las
pantallas que adaptarlas. La versión anterior se conservó en la etiqueta
\texttt{legacy/modulo-profesor-v1} y se reutilizó de ella la lógica de dibujo sobre el
documento con coordenadas relativas (\cref{SEC:DES-ANNOTATION}).

\subsection{Del contrato a los requisitos}

Para pasar del conocimiento del servicio a requisitos de interfaz se siguió un procedimiento
repetible:
\begin{enumerate}
  \item Agrupar las operaciones por etiqueta y por actor (gestor, coordinación, docente,
        estudiante).
  \item Para cada actor, ordenarlas según el ciclo de vida del recurso principal, la
        \dfn{instancia} de examen.
  \item Asociar cada transición o acción a un control con su permiso \texttt{can\_*} y su
        error esperado.
  \item Tratar cada lista paginada como una tabla con estados vacío, de carga, de error y
        sin permiso.
  \item Convertir cada condición de visibilidad del servicio en una variante de vista (por
        ejemplo, el estudiante solo ve su examen corregido con la revisión abierta).
  \item Registrar como hueco del contrato cada dato que la interfaz necesita y ninguna
        operación proporciona (\cref{SEC:AN-GAPS}).
\end{enumerate}
El procedimiento no elimina el juicio: dos decisiones no se pudieron deducir del contrato y
se tomaron por convención (la nota de aprobado, 5 sobre 10, no existe en el modelo de
datos) o tras la corrección del autor (la dirección visual, \cref{SEC:DES-VISUAL}).
